\newpage
\begin{center}
	\textbf{\large 2. Пользовательское описание}
\end{center}
% \refstepcounter{chapter}
\addcontentsline{toc}{chapter}{2. Пользовательское описание}

\section*{Обзор приложения}

Консольное приложение предназначено для локализации QR-кода на входном изображении.
В результате работы формируется JSON-файл в формате, совместимом с разметкой LabelMe: в нём сохраняются метаданные изображения и (при успешной локализации) один многоугольник, соответствующий области QR-кода.
Полученный JSON может использоваться для:
\begin{itemize}
	\item сравнения результата алгоритма с эталонной ручной разметкой (Ground Truth);
	\item визуальной проверки корректности локализации в инструментах, поддерживающих формат LabelMe;
	\item последующей передачи вырезанного фрагмента на этап декодирования QR-кода.
\end{itemize}

\paragraph{Общий формат запуска:}
\begin{verbatim}
qr_localizer <image_path> --out <out.json> [--show] [--vis <out.png>]
\end{verbatim}

\paragraph{Обязательные аргументы:}
\begin{itemize}
	\item \verb|<image_path>| -- путь ко входному изображению.
	\item \verb|--out <out.json>| -- путь к выходному JSON-файлу.
\end{itemize}

\paragraph{Опциональные флаги:}
\begin{itemize}
	\item \verb|--show| -- отображение окна с изображением и наложенным многоугольником области QR-кода.
	      Окно создаётся с параметром \verb|cv::WINDOW_NORMAL|, что позволяет изменять его размер.
	\item \verb|--vis <out.png>| -- сохранение изображения с отрисованным результатом локализации в указанный файл (например, \verb|./out/vis.png|).
\end{itemize}
